% Generated draft: requires mathematical and source review.
\documentclass[11pt]{article}
\usepackage[margin=25mm]{geometry}
\usepackage{fontspec}
\setmainfont{Latin Modern Roman}
\usepackage{amsmath,amssymb,mathtools,unicode-math}
\setmathfont{Latin Modern Math}
\usepackage{enumitem,needspace,xurl,hyperref}
\setlength{\parindent}{0pt}
\setlength{\parskip}{5pt}
\setlength{\emergencystretch}{3em}
\newcommand{\N}{\mathbb{N}}
\newcommand{\Z}{\mathbb{Z}}
\newcommand{\Q}{\mathbb{Q}}
\newcommand{\R}{\mathbb{R}}
\newcommand{\C}{\mathbb{C}}
\newcommand{\F}{\mathbb{F}}
\newcommand{\A}{\mathbb{A}}
\newcommand{\PP}{\mathbb{P}}
\newcommand{\E}{\mathbb{E}}
\newcommand{\eps}{\varepsilon}
\newcommand{\dd}{\,\mathrm{d}}
\begin{document}
\section*{UnsolvedMath ID: 30001525}
\textbf{Source problem code:} catalogue entry
\subsection{Definitions and mathematical statement}
For an integer \(n\geq 0\), let \(S_n\) denote the symmetric group on \(n\) letters, and let
\[
\Z_{(2)}=\{a/b\in\Q: a,b\in\Z,\ b\text{ odd}\}.
\]
A conventional interpretation of the \(2\)-local integral cohomology of \(S_n\) is \(H^*(S_n;\Z_{(2)})\), with trivial coefficient action. The supplied record does not explicitly specify these coefficients or whether the question concerns every finite \(S_n\), a stable cohomology object, or a collection of groups equipped with additional structure.

For clarity about the terminology, the usual coefficient sequence
\[
0\longrightarrow\Z_{(2)}\xrightarrow{\,\times 2\,}\Z_{(2)}\longrightarrow\F_2\longrightarrow 0
\]
induces a connecting homomorphism
\[
\beta:H^q(S_n;\F_2)\longrightarrow H^{q+1}(S_n;\Z_{(2)}).
\]
The associated Bockstein spectral sequence has first page \(H^*(S_n;\F_2)\); its higher differentials are commonly called higher Bocksteins. A higher-page differential is defined on a class surviving to that page, and its value is a page class. Identifying such a value with a specified integral cohomology class requires additional lifting or representative conventions. The supplied record does not specify those conventions.

\textbf{Source question.} Does the \(2\)-local cohomology of symmetric groups have an additive basis consisting of integral skyline cycles together with higher Bocksteins of skyline-basis classes?

\textbf{Unresolved qualifications.} The supplied record provides no definition of the skyline basis, of an integral skyline cycle, or of the admissible skyline classes and Bockstein orders. Consequently, it does not determine the proposed family of additive generators. Furthermore, an additive basis for a cohomology group with torsion cannot automatically mean a basis of a free \(\Z_{(2)}\)-module. If a direct-sum decomposition into cyclic summands is intended, the orders of the proposed generators and the meaning of unique expression must be specified. These missing data prevent an unambiguous quantified basis assertion. The question above therefore preserves the supplied conjectural target while leaving its source-dependent definitions and quantifiers unresolved; it does not assert the existence of such a basis.
\subsection{Short English statement}
The source asks whether integral skyline cycles and higher Bockstein classes give a complete additive description of symmetric-group cohomology localized at \(2\). A precise conjecture requires the intended groups and coefficients, the skyline constructions, the Bockstein lifting conventions, and the meaning of an additive basis in the presence of torsion. This is a statement-only draft requiring human mathematical review.
\subsection{Sources}
\begin{itemize}
\item Local UnsolvedMath registry, supplied complete record for permanent ID \(30001525\), source problem code \texttt{catalogue entry}, titled \emph{Integral Skyline Bases for Symmetric-Group Cohomology}; collection: \emph{Oberwolfach Reports Open Problems}. No individual proposer is identified. The supplied proposed year is \(2010\).
\item Supplied source citation: \emph{Cohomology of Finite Groups: Interactions and Applications} (2011), Oberwolfach Reports, pp.\ 1885--1952, DOI \texttt{10.4171/owr/2010/32}. \url{https://doi.org/10.4171/owr/2010/32}. The literature-source list also labels this reference \emph{Oberwolfach Report 32, Cohomology of Finite Groups}. The supplied background describes these reports as records of mathematical workshops and their open-problem discussions.
\item Supplied literature reference: \emph{Curtis--Wellington spectral sequence work}. \url{https://scholarsbank.uoregon.edu/server/api/core/bitstreams/e22abfd9-ad53-4d10-94ad-1d40b018ecdd/content}. The record supplies no fuller bibliographic identification or pinpoint citation for this item.
\item Supplied literature assessment, checked exactly \texttt{2026-08-21}: the record labels the problem open, reports that the Bockstein splitting statement is presented as a conjecture in the report, and states that the checked subsequent work on skyline filtrations does not establish the additive-basis conjecture in full. This is supplied dated curation derived from the OWR material, not an independently verified literature assessment. Its time-sensitive status and cited claims require independent verification.
\item The registry labels the statement \texttt{corrected\_verified} and records a manual semantic, typographic, and LaTeX review describing the cleanup as meaning-preserving, self-contained, and free of unsupported source-dependent additions. These are supplied editorial qualifications; they do not resolve the formulation gaps identified here or establish mathematical correctness.
\end{itemize}
\end{document}
